% beautybook — decorative book class with coloured theorem boxes and a
% two-tone chapter-heading treatment, adapted from the CTAN/GitHub
% BeautyLaTeX/latex-template package.
\documentclass[lang=en,12pt]{beautybook}
\RequirePackage[utf8]{inputenc}
\RequirePackage[spanish,es-tabla]{babel} % TRADUCCIÓN A ESPAÑOL
\RequirePackage{times}
\RequirePackage[T1]{fontenc}
\RequirePackage{microtype}
\RequirePackage[round]{natbib}
\RequirePackage{longtable} % Para la matriz final

\definecolor{coverbgcolor}{HTML}{e0e0e0}
\definecolor{coverfgcolor}{HTML}{00205B} % Azul marino primario
\definecolor{coverbar}{HTML}{4B6B99}
\definecolor{bottomcolor}{HTML}{00133A}
\definecolor{nuanbai}{HTML}{f5f5f5}
\coverstyle={
    cover-choose=en,
}
\mathstyle={
    math-font=plain,
}

\mynewtheorem{
    defi={\textbf{Definición}}[section]{interior style={left color=ReD!8,right color=ReD!5!CyaN!50}, borderline west={1.5mm}{0mm}{ReD}},
    thm={\textbf{Teorema}}[section]{interior style={left color=CyaN!80!black!20,right color=CyaN!80!black!15!CyaN!50}, borderline west={1.5mm}{0mm}{CyaN!80!black}},
    lem={\textbf{Lema}}[section]{interior style={left color=BluE!8,right color=BluE!5!CyaN!50}, borderline west={1.5mm}{0mm}{BluE}},
    prop={\textbf{Proposición}}[section]{interior style={left color=OrangE!8,right color=OrangE!5!CyaN!50}, borderline west={1.5mm}{0mm}{OrangE}},
    exam={\textbf{Ejemplo}}[chapter]{interior style={left color=DarkGreen!8,right color=DarkGreen!5!CyaN!50}, borderline west={1.5mm}{0mm}{DarkGreen}},
}
\newtheorem*{remark}{\textbf{Nota}}

% CAJA DE CONSOLA ESTILO CMD / MINIMALISTA
\newtcolorbox{cmdbox}[1]{
    colback=black,
    colframe=black!85,
    coltext=white!95,
    arc=0mm,
    boxrule=0.5pt,
    left=3mm,
    right=3mm,
    top=1mm,
    bottom=1mm,
    before skip=1em,
    after skip=1.5em,
    title={\texttt{>\_ \textbar\ \ #1}},
    fonttitle=\ttfamily\small,
    breakable
}

\begin{document}
\thispagestyle{empty}
\title{Cybersecurity Framework 2.0 en la UNAP}
\subtitle{Análisis de Procesos de Ataque y Defensa}
\edition{}
\bookseries{}
\author{QUISPE CHATA CARLOS DAVID}
\pressname{UNAP}
\presslogo{inner_pics/logo-universidad_files/Logo_UNAP_x7pS.png}
\makecover

\makeatletter
\definecolor{bg}{HTML}{E6EEF5} % Azul claro para los fondos de secciones
\definecolor{fg}{HTML}{00205B} % Azul marino (Color primario absoluto)
\colorlet{outermarginbgcolor}{bg}
\colorlet{outermarginfgcolor}{fg}
\oddoutermargin{\sffamily \leftmark}
\evenoutermargin{\sffamily\@title}
\chapimage{inner_pics/beautybook-mark.png}
\makeatother

\colorlet{framegolden}{fg}
\colorlet{framegray}{bg!50}

\frontmatter
\pagenumbering{Roman}

\thispagestyle{empty}
\tableofcontents\let\cleardoublepage\clearpage

\mainmatter
\pagenumbering{arabic}

\partabstract{\hspace*{2em} Breve resumen o introducción a la primera parte del nuevo trabajo.} Este informe desarrolla la implementación técnica, normativa y operativa del NIST Cybersecurity Framework (CSF) 2.0 en la Universidad Nacional del Altiplano, integrando metodologías defensivas y tácticas ofensivas.}
\part*{Primera Parte}



% ==========================================
% NUEVO CONTENIDO (TRABAJO 3)
% ==========================================

\chapter{Nuevo Tema}
\label{ch:nuevo_tema}

Aquí comienza el contenido del nuevo tema. Utiliza los comandos y entornos ya definidos en el preámbulo (como \texttt{cmdbox}, etc.).


\end{document}
